\SourceLeaf{216}$cbe = cb'b$. Es wird also $be$ die 2 schneiden, in $e'$. Macht man nun $b'g = be'$, so wird $cg$ und $cd'$ mit $cb'$ einerlei Winkel machen. Welches absurd ist.

5. Wenn 1 mit 2 und 1 mit 3 parallel, so ist auch 2 mit 3 parallel.

6. Was correspondirende Punkte auf zwei Parallelen sind.

7. Aequidistanz der correspondirenden Punkte.

8. Der Punkt auf einer dritten Parallelen correspondirt correspondirenden Punkten auf den beiden ersten.

9. \textls{Trope} ist die L[inie, die von correspondirenden Punkten gebildet wird, wenn man alle Parallelen zu einer Geraden betrachtet.]

\Needspace{19\baselineskip}
\bigskip
\begin{center}
\rule{2.5cm}{0.4pt}

\textsc{Bemerkung.}
\end{center}

In dem Briefe an \textsc{Schumacher} vom 17. Mai 1831 (S. 213 dieses Bandes) sagt \textsc{Gauss}, dass er von seinen Meditationen über die Theorie der Parallellinien, die schon gegen 40 Jahr alt seien, früher nie etwas aufgeschrieben habe, dass er jedoch vor einigen Wochen begonnen habe, einiges aufzuschreiben; auch in dem Briefe an \textsc{Bolyai} vom 6. März 1832 erwähnt \textsc{Gauss} solche Aufzeichnungen. Man wird daher nicht fehlgehen, wenn man annimmt, dass die vorstehenden nicht datirten Notizen, die sich auf einzelnen Zetteln befinden, aus dem Jahre 1831 stammen.

Die Figuren in der Notiz [1, sowie die erste Figur in der Notiz [3] sind dem Texte hinzugefügt worden

\hfill \textsc{Stäckel}.

